\documentclass[11pt]{article}

\usepackage[utf8]{inputenc}
\usepackage[T1]{fontenc}
\usepackage{times}
\usepackage[textwidth=5.5in,textheight=9in]{geometry}
\usepackage{graphicx}
\usepackage{booktabs}
\usepackage{float}
\usepackage{amsmath,amsfonts}
\usepackage{xcolor}
\usepackage{hyperref}
\usepackage{todonotes}


\newcommand{\note}[1]{\textcolor{blue}{#1}}

\title{\vspace{-1em}\textbf{Trustworthiness of Models based on Language}\\
  \vspace{1em}
  \normalsize{\normalfont Responsible AI Project Report}}

\author{
  Hannah Buch \quad Yassine Ajoud\\
}

\date{}

\begin{document}

\maketitle

\begin{abstract}
This report details the development of an evaluation framework designed to assess the capability of Large Language Models (LLMs)—specifically Qwen2.5-7B-Instruct—to recognise and analyse literature across multiple languages. Using text snippets automatically extracted and processed from Project Gutenberg, we tested the model's capacity to identify the correct book title from a set of four candidates. The evaluation spanned five languages: English, German, French, Italian, and Dutch. 

Beyond standard recognition, the framework incorporated several tests: cross-lingual evaluation (matching German translations to English titles to test translation application), robustness testing via noisy text, and an analysis of the model's reasoning for its choices. All results were visualised graphically. Across all variants, English consistently performed better than the other languages, indicating that the model's literary recognition capabilities are highly language-dependent, leading to weaker performance for users operating in non-English contexts.
\end{abstract}

\section{Introduction}
The rise of Artificial Intelligence in the last years, saw especially a big development in LLMs, Large Language models. Models who were trained on large chunks of text, and capable of reproducing text, based on this trainings data. Large language models answer questions in many languages, yet most of their training data is in English. The training data mainly stems from the internet, making it in large parts english and european centric. It is therefore reasonable to ask whether they perform as well in other languages as they do in English. The question
matters regarding trust and fairness. A model should be usable for all types of users, including those with different linguistic backgrounds. When a user in Germany addresses the model in German, they would not be able to notice whether the output has a better or worse quality than in other languages. 

The suggested course project approach was to take an existing benchmark, translate it with a chatbot and compare the scores before and after. During the project, we departed from this plan. While we stayed within the realm of assessing literature knowledge of LLMs, we shifted the focus.  
While machine translations often count as reliable, at least when done on well-trained languages, this mostly extends to a translation in a more technical way; the translation will probably carry the correct meaning, but it will lack the same style and artistic approach. Human book translations are well known for translating more the style of the book than the single words. So for us, the question became: how well can an LLM work with translations that cannot just be translated word for word back into the original? With such a base of books at our disposal, we used literature from the Gutenberg Project, where we found works from England, Germany, France, Italy, and the Netherlands. Project Gutenberg is a large online library consisting of books, which are not under any copyright, mostly due to the age of the text. It is reasonable to assume that the most current LLMs were trained on this library as it is easy to access and has most of the well-known classics in its database. 

For the project we took the model through different types of tests, assessing its capabilites. The focus was to test reproducability of training data, but with some obstacles, making it more difficult. The main approach was to test the model on text parts and asking for a decision on a multiple-choice basis. A model that encountered the book during training should do this well, while a model that never saw the book can still guess from the names and places in the passage. Random guessing would yield 25\%.

We ran this test with Qwen2.5-7B-Instruct on 100 books each in English,
German, French, Italian and Dutch. On top of these five runs we carried out
three smaller experiments. In the first, we took 23 German translations of
English books and asked the model for the English original title. In the
second, we replaced characters in the passages with random letters
and observed the effect on accuracy. In the third, we required the model to
explain its answer on smaller sets of books in English, German, and French.

\section{Related Work}
Chang et al. investigated which books ChatGPT and GPT-4 have seen between books in and out of the copyright\cite{chang2023}. They took short passages from novels, masked a single character name and let the model fill in the missing name. When a model supplies the correct name repeatedly, it has most likely encountered the book during training. They found that the models have memorised a large number of books and that the books whose passages appear most often on the web are memorised best. Models were tested on copyrighted and free books, where especially the open books had good results. This paper implied a large background knowledge of LLMs regarding well-known literature, which was an important starting point for our testing, assuming that the model had seen the books in training. Our setup also used a similar experiment approach, but focusing on titles instead. We also extended the testing towards another model,  as most paper used ChatGPT as base model.
Zhang et al. studied the gap between English and other languages in large language models \cite{zhang2023}. They gave the same tasks to the models in several languages and compared the results.For some tasks the correct answer stays the same after translation, while for others it depends on the language itself. On these language-dependent tasks the  models had clear difficulty once the question was no longer in English. The paper comes to the conclusion, that english functions as the models native language. Eventhough reasoning and knowledge should supposedly be independent from languages, english had consistently higher accuracies. The paper also speculates whether the models internally translate inputs into english. This would be a problem for the models in our tests, as books are usually not translated word for word, therefore this approach to literature analysis would be more complicated for the models. Also, because titles in different languages are not necessarily just translations of each other, but at times, something completely different. 
But not all paper were critical towards the results of LLMs in different languages. Shi et al. had a more positive take on the language results, regarding multi-language reasoning\cite{shi2022}. They tested chain of thought reasoning in different languages, and hat really good results, even in languages that make up a very small percentage of the training corpus. In the end of their paper they also tested how well models could understand the sense of words across different languages, where they had better results, than previous research when applying chain of thought. But they still concluded that this is a difficult task. Something we later apply towards literature translations, as human translations usually translate the sentiment but not necessarily the words, of a text.
Despite language tests in models being done in large quantities, there still seems to be a high dependence of the results on the task given. While training a model on given good reasoning, seemed to give good results, we now wanted to test, when the model was only trained on the texts, but not the task it was supposed to perform. Based on prior research and our own tests with the model, we performed tasks similar to the related works, but implementing different kinds of obstacles for the model, instead of a straighforward approach. 
\section{Approach}
\subsection{Question and Goals}
The primary objective we intended to investigate was: How accurate are LLMs at performing zero-shot literary identification when input and output are supposed to be different languages. While models show strong performance in one singular language, we wanted to extend this testing to cross-lingual. While working on this we developed an additional question.  How resilient are LLMs capabilites when tasked with identifying corrupted literary text? Old digitized archives, but also old grammar or regional dialects, can lead to errors or intentional differences in text, which possibly hinder the recognition process. Therefore we tried to classify how much corruption still yields acceptable results. 
To answer this questions we built an automated benchmarking pipeline, using data streamed from Project Gutenberg. The results are visualized via graphs for later analysis. To ensure reproducibility and interactive execution we wrote the program into a Jupyter Notebook, which not only made the graphics easier to display but also enables smaller checks, without running the whole program. 

\subsection{Model and Database}
For the database, we searched through Hugging Face for appropriate literature databases to use. First, we wanted to use nikolina-p/gutenberg\_clean\_en\_splits\footnote{https://huggingface.co/datasets/nikolina-p/gutenberg\_clean\_en\_splits}, because the dataset already contained only books in English and had stripped all book headers. At first, this seemed like a good approach to make sure the model would only be asked about actual text parts, and not accidentally some part of the header which did not contribute to the book itself. Unfortunately, the books in this database were only identifiable by the numbers anymore, as title and similar things were contained in the header. That is why we switched to manu/project\_gutenberg\footnote{https://huggingface.co/datasets/manu/project\_gutenberg/viewer/default/en}, which was the original database used by nikolina-p/gutenberg\_clean\_en\_splits. This dataset holds the raw text of Project Gutenberg books, with additional information in the header text. The dataset was streamed, so that the code reads one book at a time over the network.
The model is Qwen2.5-7B-Instruct \cite{qwen2025}, an open instruction tuned model with 7 billion parameters that we run on a local machine through the transformers library. It was especially chosen for its known good performance, and its size of parameters. This size of model is still reasonable to execute on a local computer. 

\section{Experiments}

Table~\ref{tab:overview} lists all runs. Each run used the same model, the
same greedy decoding and a single pass over its question set with no
repetition. Every number in the next section therefore comes from one pass.
The five base runs use 100 books per language, while the three additional runs
are smaller.

\begin{table}[H]
\centering
\small
\begin{tabular}{lllll}
\toprule
Run & Text & Titles & Books & Snippet \\
\midrule
English & English & English & 100 & 1000 chars, middle \\
German & German & German & 100 & 1000 chars, middle \\
French & French & French & 100 & 1000 chars, middle \\
Italian & Italian & Italian & 100 & 1000 chars, middle \\
Dutch & Dutch & Dutch & 100 & 1000 chars, middle \\
Translations & German & English & 23 & 3000 chars, early \\
Noise & German & German & 20 & 1000 chars, 50\% noise \\
Explanations & English & English & 10 & 1000 chars, middle \\
\bottomrule
\end{tabular}
\caption{Overview of all runs.}
\label{tab:overview}
\end{table}

\subsection{Results}

Table~\ref{tab:results} reports the outcome of all runs. Every base run stays
well above the 25\% guess line. English is the strongest set with 88 of 100
books, followed by Dutch with 83. German and French sit close together at 72
and 71, and Italian is the weakest at 66.

\begin{table}[H]
\centering
\small
\begin{tabular}{lrrr}
\toprule
Test set & Books & Correct & Accuracy \\
\midrule
English & 100 & 88 & 88.0\% \\
Dutch & 100 & 83 & 83.0\% \\
German & 100 & 72 & 72.0\% \\
French & 100 & 71 & 71.0\% \\
Italian & 100 & 66 & 66.0\% \\
\midrule
German text, English titles & 23 & 22 & 95.7\% \\
German with 50\% noise & 20 & 6 & 30.0\% \\
English with explanations & 10 & 5 & 50.0\% \\
\bottomrule
\end{tabular}
\caption{Accuracy of Qwen2.5-7B-Instruct on all runs. Random guessing gives
25\%.}
\label{tab:results}
\end{table}

The translation run reached 22 of 23 books. The model read German text and
still selected the correct English original title in almost every case.

The noise run fell sharply. Without noise the German base run stands at 72\%,
but once half of the characters are replaced the model got 6 of 20 books
right. That is 30\% against a guess line of 25\%, so at this noise level the
recognition ability is largely gone.

The explanation run reached 5 of 10 books where the base English run answered
6 of the same 10 questions correctly. The samples are tiny, so a difference of
one book carries little weight. The content of the explanations is more
revealing than the score. For book 1 the model chose the right option and
named \textit{The Tatler} as the clue in the text. For book 2 it chose a wrong
option and wrote that the passage is clearly from \textit{The Prodigals and
their Inheritance} because the character Mr.Probert is specific to that book.
Mr.Probert is in fact a character in \textit{The Reverberator}, which was the
correct answer. The model sounded equally certain in both cases.

\section{Analysis}

The plain reading of Table~\ref{tab:results} is that the model recognizes
literature best in English. German scores 16 points below English and French
17, while Italian is 22 points below. Dutch loses only 5 points, which is higher
than we expected. Part of the reason may be a difference in filtering. The
English and German runs removed translated books before building questions,
whereas the French, Italian and Dutch runs kept them because the translation
filter was not active there. A Dutch or French copy of a famous English book
is easier to recognize than a local original so the gap for these three
recognise is probably somewhat larger than the table suggests.

The 95.7\% for German translations with English titles looks very high next to
the 72\% for German books with German titles and we would not read too much
into it. The set is small at 23 books and the books are famous  with Dickens,
Twain, Wilde, Shakespeare and Carroll all present. Eleven files are volumes of
the same history so the same correct title recurs eleven times. The passages
for this run were also three times longer and the answer pool was small. All
of this pushes the score upward.
 The model can link German text to an English title when the book is well known,
which fits the finding of Chang et al. \cite{chang2023}. Famous books have
many copies on the web in several languages so the model had many
opportunities to learn the link during training.

The noise run gives a clear picture. At 50\% random characters the score lands
at 30\% against a guess line of 25\%. A passage in which every second
character is random is hard for a human to read as well. A single noise level
is of course a thin basis.

The explanation run is small yet it exposes a known problem from a new angle.
The model states wrong facts with full confidence. The Mr.Probert case
is a good example, as the explanation picks a real detail from the
passage and attaches it to the wrong book. A user who does not know both books
cannot catch the error. For trustworthiness this weighs more than the raw
accuracy numbers because a wrong letter is cheap while a wrong letter wrapped
in a fluent and confident account can mislead a reader.

Some limitations remain. Every run is a single pass so the numbers carry
sampling luck and with 100 books a swing of a few points between repeated
runs would be normal. The titles come from the file headers and were
lowercased, so a few odd or broken titles may have entered the question sets.
The three wrong titles per question are random draws, which makes some
questions easy when the wrong titles fall far from the topic of the passage.

\section{Conclusion}

We built a small test around one clear question. Does a language model
recognize books in other languages as well as it does in Enrecogniser
Qwen2.5-7B-Instruct the answer is no. English reached 88\% while the other
four languages landed between 66\% and 83\%. The model beats random guessing
everywhere, so the ability does not vanish outside English but it does weaken.
The additional runs sharpen the picture. The model links German translations
to their English originals when the books are famous, it breaks down when half
of the text is random noise and its explanations sound convincing even when
they are wrong.

\section*{Author contribution statement}
Both authors contributed to the project and discussed the
design and code implementation. Hannah Buch was the main contributor in building
and running the experiments. Yassine Ajoud focused on the written report
and the presentation.


\section*{GenAI use}
DeepL was used for simple translations and for wording improvements in
the written report.

\begin{thebibliography}{9}

\bibitem{chang2023}
Kent K. Chang, Mackenzie Cramer, Sandeep Soni and David Bamman.
Speak, Memory: An Archaeology of Books Known to ChatGPT/GPT-4.
In \textit{Proceedings of the 2023 Conference on Empirical Methods in
Natural Language Processing (EMNLP)}, 2023. arXiv:2305.00118.

\bibitem{zhang2023}
Xiang Zhang, Senyu Li, Bradley Hauer, Ning Shi and Grzegorz Kondrak.
Don't Trust ChatGPT when Your Question is not in English: A Study of
Multilingual Abilities and Types of LLMs.
In \textit{Proceedings of the 2023 Conference on Empirical Methods in
Natural Language Processing (EMNLP)}, 2023. arXiv:2305.16339.

\bibitem{shi2022}
Freda Shi, Mirac Suzgun, Markus Freitag, Xuezhi Wang, Suraj Srivats, Soroush Vosoughi, Hyung Chung, Yi Tay, Sebastian Ruder, Denny Zhou, Dipanjan Das and Jason Wei.
Language Models are Multilingual Chain-of-Thought Reasoners.
In \textit{arXiv preprint arXiv:2210.03057}, 2022.


\end{thebibliography}

\end{document}
